\section{Prediction intervals, detection, and the sizing problem}
\label{sec:sizing}

An identifiable exponent is an input to capacity planning. This section
assembles a sizing framework: a throughput chance constraint and an optional
detection constraint. Any detector that yields measured $(\pd,\pf,\rho)$ plugs
in; experiment~X5 is one instance.

\subsection{Definitions and assumptions}

\begin{Def}[Fault class and detector]
\label{def:faults}
A subset $B$ of size $\kf$ is \emph{faulty} if each member fails a proposal
round independently with probability $\pd>\pf$ (omission class; equivocation
excluded). The detector flags validator~$i$ when its participation over $\w$
rounds falls below $\thr\in(\pf,\pd)$; detection succeeds when at least one
member of~$B$ is flagged.
\end{Def}

\begin{Assumption}[Exchangeable correlation]
\label{as:rho}
Participation indicators in~$B$ are exchangeable with pairwise correlation
$\rho\in[0,1)$, so the mean carries design effect $1+(\kf-1)\rho$ in the sense
of Kish~\cite{kish1965}.
\end{Assumption}

\subsection{Prediction intervals and detection bound}

With regressor $\mathbf{z}=(1,\log\cc,-\log\n)^{\!\top}$ and
$\hat\vartheta=(\log\Tzero,\gamma,\beta)^{\!\top}$, the delta
method~\cite{seber2003} gives predictive variance
$\hat{V}=\mathbf{z}^{\!\top}\widehat{\Cov}(\hat\vartheta)\mathbf{z}+\hat\sigma^{2}$.

\begin{State}[Lognormal prediction interval]
\label{st:pi}
An approximate $(1-\epsp)$ interval for $\TPS(\n,\cc)$ is
$\exp\{\log\widehat{\TPS}\mp\zq_{1-\epsp/2}\hat{V}^{1/2}\}$; the one-sided lower
endpoint $\underline{\TPS}_{1-\epsp}$ uses $\zq_{1-\epsp}$.
\end{State}

A residual bootstrap~\cite{efron1993} supplies empirical quantiles when residual
normality is doubtful. Coverage is checked on held-out validator counts
(Section~\ref{sec:exp}).

\begin{State}[Approximate miss probability and $\nmin$]
\label{st:rho}
With $\kf=\lceil\phi\n\rceil$, the design-effect heuristic yields
\begin{equation}
\pmiss(\n,\w)\lesssim\exp\!\Bigl\{-\frac{2\w\kf(\pd-\thr)^{2}}{1+(\kf-1)\rho}\Bigr\},
\label{eq:pmiss}
\end{equation}
which is nonincreasing in $\n$ and $\w$. Treating~\eqref{eq:pmiss} as an
engineering constraint $\pmiss\le\epss$ gives the cut
\begin{equation}
\nmin=\left\lceil
\frac{1+(\kf-1)\rho}{\phi}\cdot
\frac{\log(1/\epss)}{2\w(\pd-\thr)^{2}}\right\rceil .
\label{eq:nmin}
\end{equation}
\end{State}

\begin{HeuristicJustification}
A miss requires every member of~$B$ to remain unflagged. Classical
Hoeffding~\cite{hoeffding1963} applies to independent bounded summands.
Following the survey-sampling design-effect approximation~\cite{kish1965}, we
replace the correlated sample by an effective size
$\neff=\kf\w/(1+(\kf-1)\rho)$ and insert that size into the independent
Hoeffding form. Expression~\eqref{eq:pmiss} is thus an engineering adaptation
for dependent trials; $\pd$, $\pf$ and $\rho$ are measured in experiment~X5.
\end{HeuristicJustification}

The detection model is an architectural framework slot: it accepts future
high-fidelity metrics $(\pd,\pf,\rho)$ without changing the optimization. With
an informative limit of AUC$\,=\,\resAuc$ under encrypted transport, the sizing
framework defaults to the throughput-bound constraint and the optimization
window remains well-defined.

\subsection{The sizing program}

\begin{Problem}[Validator-set sizing]
\label{prob:sizing}
Given peak demand $\Dpk$, levels $\epsp,\epss$, window~$\w$, faulty fraction
$\phi$, production scale-up $\kappaScale>0$ from \RNM{} units to absolute
throughput, and nondecreasing $\Cost(\n)$, minimize $\Cost(\n)$ subject to
$\Prob[\kappaScale\cdot\TPSn(\n,\cc^{\star})\ge\Dpk]\ge1-\epsp$ and
$\pmiss(\n,\w)\le\epss$, where $\cc^{\star}=\arg\max\gfun$. When $\kappaScale$
is free, take the smallest value that meets the throughput constraint at
$\n=4$.
\end{Problem}

The chance constraint is equivalent, via Statement~\ref{st:pi}, to
\begin{equation}
h(\n)\eqdef\log(\kappaScale\,\Tzero\,\gfun(\cc^{\star})/\Dpk)
-\beta\log\n-\zq_{1-\epsp}\hat{V}^{1/2}(\n,\cc^{\star})\ge0.
\label{eq:h}
\end{equation}
On the extrapolation range the fitted covariance confirms that $\hat{V}$ is
nondecreasing in $\log\n$ (Section~\ref{sec:exp}), so with $\beta>0$ the map
$h$ is strictly decreasing there; on the measured grid we evaluate~$h$
directly.

\begin{Proposition}[Feasibility window]
\label{prop:window}
Whenever $h$ is strictly decreasing and the approximate detection cut
of Statement~\ref{st:rho} is used, the feasible set is
$\{\n:\nmin\le\n\le\nmax\}$ with $\nmin$ from~\eqref{eq:nmin} and
$\nmax=\lfloor\sup\{\n:h(\n)\ge0\}\rfloor$. If nonempty and $\Cost$ is
nondecreasing then $\nopt=\nmin$.
\end{Proposition}

\begin{Proof}
Detection holds on $\{\n\ge\nmin\}$ and throughput on $\{\n\le\nmax\}$, hence an
interval; monotone cost selects the left endpoint.
\end{Proof}

\begin{Corollary}[Empty feasible set]
\label{cor:infeasible}
If $\nmin>\nmax$, no validator count satisfies both requirements: the
architecture must change, not merely the provisioning.
\end{Corollary}

The window is the intersection of two monotone constraints
\cite{charnes1959,nemirovski2007}; it converts calibrated exponents into a
sizing rule.
